% Cuerpo compartido. Rutas relativas a la carpeta desde la que se compila
% (master/ o la carpeta de cada alumno, ambas a un nivel de documentation/).
% Cap. 1-5: Seminario y TFI. Cap. 6-9: solo TFI (\tfitrue).
% Cap. 10 y Anexos: propios de cada carpeta (casos de uso distintos por alumno).

\newcommand{\capitulo}[1]{\capredactadofalse\input{#1}}

\part*{Plan de Negocios}
\addcontentsline{toc}{part}{Plan de Negocios}
\capitulo{../master/01_descripcion_general/texto}
\capitulo{../master/02_analisis_estrategico/texto}
\capitulo{../master/03_analisis_foda/texto}
\capitulo{../master/04_segmentacion/texto}
\capitulo{../master/05_plan_de_accion/texto}
\iftfi
\capitulo{../master/06_plan_de_marketing/texto}
\capitulo{../master/07_operaciones/texto}
\capitulo{../master/08_plan_financiero_economico/texto}
\capitulo{../master/09_factibilidades/texto}
\fi

\ifparcial\else
% La guía indica que la carpeta técnica debe presentarse como cuerpo separado.
\part*{Carpeta Técnica}
\addcontentsline{toc}{part}{Carpeta Técnica}
% El Seminario omite los Cap. 6-9, pero la carpeta técnica es siempre el Cap. 10.
\iftfi\else\setcounter{section}{9}\fi
% El Cap. 10 es individual: cada carpeta (y master/, como plantilla) tiene su propio texto.tex.
\capitulo{10_solucion_tecnologica/texto}
\capitulo{../master/11_contenido/texto}

\section*{Anexos}
\addcontentsline{toc}{section}{Anexos}
Información ampliatoria (CV, folletería, informes de investigación, normas, textos de convenios).
Evitar mezclar este material con el cuerpo principal del documento.

\fi

\printbibliography[title={Referencias}]
